\appendix

\appendixpage
\addappheadtotoc

\section{Modelo fenomenológico do sistema reator-separador}
\label{sec:apendice_modelo}

O modelo fenomenológico completo da rede reator-separador, formulado por \textcite{pourkargar_2017}, é constituído por doze equações diferenciais, deduzidas a partir de balanços de massa e de energia nos reatores e no separador:
\begin{align}
  \frac{dV_1}{dt}    & = F_{f1} + F_R - F_1                                                                                                                                               \\
  \frac{dV_2}{dt}    & = F_{f2} + F_1 - F_2                                                                                                                                               \\
  \frac{dV_3}{dt}    & = F_2 - F_P - F_R - F_3                                                                                                                                            \\
  \frac{dT_1}{dt}    & = \frac{F_{f1}}{V_1}(T_0 - T_1) + \frac{F_R}{V_1}(T_3 - T_1) + \frac{Q_1}{\rho C_p V_1} - \frac{m}{C_p}(k_{11} x_{A1} \Delta H_{r1} + k_{21} x_{B1} \Delta H_{r2}) \\
  \frac{dT_2}{dt}    & = \frac{F_{f2}}{V_2}(T_0 - T_2) + \frac{F_1}{V_2}(T_1 - T_2) + \frac{Q_2}{\rho C_p V_2} - \frac{m}{C_p}(k_{12} x_{A2} \Delta H_{r1} + k_{22} x_{B2} \Delta H_{r2}) \\
  \frac{dT_3}{dt}    & = \frac{F_2}{V_3}(T_2 - T_3) + \frac{Q_3}{\rho C_p V_3}                                                                                                            \\
  \frac{dx_{A1}}{dt} & = \frac{F_{f1}}{V_1}(x_{A0} - x_{A1}) + \frac{F_R}{V_1}(x_{AR} - x_{A1}) - k_{11} x_{A1}                                                                           \\
  \frac{dx_{B1}}{dt} & = \frac{F_R}{V_1}(x_{BR} - x_{B1}) - \frac{F_{f1}}{V_1} x_{B1} + k_{11} x_{A1} - k_{21} x_{B1}                                                                     \\
  \frac{dx_{A2}}{dt} & = \frac{F_{f2}}{V_2}(x_{A0} - x_{A2}) + \frac{F_1}{V_2}(x_{A1} - x_{A2}) - k_{12} x_{A2}                                                                           \\
  \frac{dx_{B2}}{dt} & = \frac{F_1}{V_2}(x_{B1} - x_{B2}) - \frac{F_{f2}}{V_2} x_{B2} + k_{12} x_{A2} - k_{22} x_{B2}                                                                     \\
  \frac{dx_{A3}}{dt} & = \frac{F_2}{V_3}(x_{A2} - x_{A3}) - \frac{F_P + F_R}{V_3}(x_{AR} - x_{A3})                                                                                        \\
  \frac{dx_{B3}}{dt} & = \frac{F_2}{V_3}(x_{B2} - x_{B3}) - \frac{F_P + F_R}{V_3}(x_{BR} - x_{B3})
\end{align}
em que o índice $i$ identifica a unidade do processo: $i=1$ corresponde ao primeiro reator, $i=2$ ao segundo reator e $i=3$ ao separador.
Os termos $V_i$ e $T_i$ denotam, respectivamente, o volume e a temperatura de cada unidade, e $T_0$ é a temperatura da corrente de alimentação.
As vazões volumétricas de alimentação dos reatores são $F_{f1}$ e $F_{f2}$; $F_1$ e $F_2$ são as vazões que deixam os reatores 1 e 2, $F_3$ é a vazão de produto retirada do separador, $F_R$ é a vazão de reciclo e $F_P$, a de purga.
A taxa de calor fornecida a cada unidade é $Q_i$, ao passo que $\rho$ e $C_p$ são a massa específica e o calor específico do fluido, $m$ é a molalidade e $\Delta H_{rj}$ é a entalpia da reação $j$.
A composição do meio é dada pelas frações molares $x_{Ai}$ e $x_{Bi}$ dos componentes $A$ e $B$ na unidade $i$.
Por fim, $x_{A0}$ é a fração molar de $A$ na alimentação, $x_{AR}$ e $x_{BR}$ são as frações molares de $A$ e $B$ na corrente de reciclo e $k_{ri}$ é o coeficiente de velocidade da reação $r$ na unidade $i$.

Os coeficientes de velocidade das reações, $k_{ri}$, são funções da temperatura da unidade $i$ e seguem a lei de Arrhenius.
Os valores dos parâmetros são os reportados por \textcite{pourkargar_2017}.

\clearpage

\section{Ponto de Operação do Modelo}
\label{sec:apendice_ponto_operacao}

O ponto de operação em regime estacionário adotado para a linearização do sistema é o reportado por \textcite{pourkargar_2017}.
Os valores das entradas e dos estados nesse ponto constam, respectivamente, nas Tabelas~\ref{tab:u_ss} e \ref{tab:y_ss}.

\begin{table}[h]
  \caption{Valores das variáveis de entrada no ponto de operação em regime estacionário ($u_{ss}$).}
  \label{tab:u_ss}
  \centering
  \begin{tabular}{lll}
    \hline
    Variável & Valor   & Unidade \\
    \hline
    $F_{f1}$ & 5{,}04  & m³/h    \\
    $F_{f2}$ & 5{,}04  & m³/h    \\
    $F_R$    & 17{,}00 & m³/h    \\
    \hline
  \end{tabular}
\end{table}


\begin{table}[h]
  \caption{Valores dos estados do sistema no ponto de operação em regime estacionário ($y_{ss}$).}
  \label{tab:y_ss}
  \centering
  \begin{tabular}{lll}
    \hline
    Variável & Valor   & Unidade \\
    \hline
    $T_1$    & 432{,}4 & K       \\
    $T_2$    & 427{,}1 & K       \\
    $T_3$    & 432{,}1 & K       \\
    $x_{A1}$ & 0{,}536 & --      \\
    $x_{B1}$ & 0{,}448 & --      \\
    $x_{A2}$ & 0{,}545 & --      \\
    $x_{B2}$ & 0{,}438 & --      \\
    $x_{A3}$ & 0{,}298 & --      \\
    $x_{B3}$ & 0{,}670 & --      \\
    \hline
  \end{tabular}
\end{table}


\clearpage
